\section{What adaptive evaluation costs}
\label{sec:cost}

Adaptive testing saves items only after somebody has answered the whole bank: the
calibration panel. Counting one administered item as the unit of cost, evaluating $M$ new
systems costs $MI$ with the full bank of $I$ items, and $NI+MT$ with a fresh panel of $N$
calibrators and adaptive tests of $T$ items, so adaptive testing pays off once
$M>NI/(I-T)$. With eight calibrators and $T=100$ this break-even lies between
\CostBreakFullMin\ and \CostBreakFullMax\ new systems on every benchmark
(Table~\ref{tab:cost}), and for
100 new systems the total cost is about \CostRatioHundred\% of full evaluation.

The fairer comparison is with uniform subsampling, which needs no panel at all. At
matched precision, EOST-M needed a median of \CostItemsNaive\ items to reach the score
error that 100 uniformly sampled items reach with the sample mean, so it saves about
\CostSaveNaive\ items per system but must first pay $NI$; the break-even is then a median of
\CostBreakNaiveMed\ new systems (from \CostBreakNaiveMin\ to \CostBreakNaiveMax\ across
benchmarks). Against Random:mg, which pays for the same panel, adaptivity alone saves
items from the first system on \CostRandomSaveBanks\ of 8 benchmarks: EOST-M matched 100
random items scored by the same model after a median of \CostItemsRandom\ items. Two situations change the arithmetic. When the
calibrators' full results already exist---public leaderboards, previous internal
evaluations---the panel is free and adaptive testing is cheaper from the first system.
When items are expensive to administer and a benchmark is used for many systems, as for
agentic benchmarks, the per-system saving dominates quickly. The computation itself is
negligible: calibrating a panel and running a 100-item test takes seconds on one CPU core
(Section~\ref{sec:design}).

\begin{table}[tbp]
\centering\footnotesize
\caption{Cost of adaptive evaluation with a fresh calibration panel of $N=8$ systems that must first answer all $I$ items (study 1). Break-even: number of new systems from which EOST-M with 100 items costs less than full evaluation; items for which EOST-M matches the score error of 100 uniformly sampled items scored by the sample mean (Naive) or by the same calibrated model (Random:mg), and the resulting break-even against Naive. When the calibrators' full results already exist (public leaderboards), the calibration cost vanishes and EOST-M is cheaper from the first system.}
\label{tab:cost}
\begin{tabular}{@{}lrrrrr@{}}
\toprule
 & & break-even vs & items to match & break-even vs & items to match \\
Benchmark & $I$ & full evaluation & Naive@100 & Naive & Random:mg@100 \\
\midrule
GSM8K & 1319 & 8.7 & 70 & 352 & 80 \\
ARC-C & 1172 & 8.7 & 20 & 117 & 90 \\
HellaSwag & 2000 & 8.4 & 8 & 174 & 25 \\
WinoGrande & 1267 & 8.7 & 20 & 127 & 60 \\
MMLU & 2000 & 8.4 & 15 & 188 & 100 \\
SWE-bench V. & 500 & 10.0 & 40 & 67 & 70 \\
LiveCodeBench & 1055 & 8.8 & 25 & 113 & 60 \\
CRUXEval-O & 800 & 9.1 & 25 & 85 & 60 \\
\bottomrule
\end{tabular}
\end{table}

